\input{preamble.tex}
\textbf{\large Assignment 10: Best Practices in Computing}\\\vspace{10pt}
\end{center}

\vspace{10pt}
\noindent
\textbf{\large Instructions:}

\vspace{10pt}
\begin{itemize}
\setlength\itemsep{0pt}
  \item {\color{red}{\textbf{Deadline}}}: \textbf{April 30, before class}
  \item Submit via your GitHub repository in a separate folder (\texttt{assignment10})
  \item (No solutions this time, but no take home part either)
\end{itemize}

\vspace{20pt}
\tableofcontents
\newpage

% ==========================================================================
\section{Task: Build a Mini Research Project}
% ==========================================================================

The goal of this assignment is to build a small, self-contained project that integrates R analysis with a \LaTeX{} document. The idea is simple: R produces output files (a table and a figure), and the \LaTeX{} document loads them automatically. When the analysis changes, you recompile the document and it updates --- no copy-pasting.

{\sloppy Use any dataset from previous assignments, available at the course GitHub repository (\href{https://github.com/franvillamil/AQM2/tree/master/datasets/}{github.com/franvillamil/AQM2/tree/master/datasets/}).\par}

\subsection{What to build}

Create a folder \texttt{assignment10} in your course repository with the following components:

\begin{enumerate}[label=\arabic*.]

  \item \textbf{An R script} (e.g., \code{analysis.R}) that:
  \begin{itemize}
    \item Loads the dataset
    \item Runs at least two regression models (e.g., CPI on GDP with and without controls, using the \code{corruption} dataset)
    \item Saves a regression table as a \code{.tex} file using \code{modelsummary()} with the \code{output} argument
    \item Creates at least one figure (e.g., a scatter plot) and saves it as a \code{.pdf} file using \code{ggsave()}
  \end{itemize}

  \item \textbf{A \LaTeX{} document} (e.g., \code{document.tex}) that:
  \begin{itemize}
    \item Includes a brief title and one or two sentences describing the analysis
    \item Loads the regression table using \verb|\input{path/to/table.tex}|
    \item Loads the figure using \verb|\includegraphics{path/to/figure.pdf}|
    \item Compiles into a PDF that shows the table and figure
  \end{itemize}

\end{enumerate}

The specifics of the analysis, the model specification, the plot design, and the folder structure are up to you. What matters is that the pipeline works: run the R script, compile the \LaTeX{} document, and you get a PDF with the results embedded automatically.

\textbf{Note:} at the very least, I recommend creating directories \code{img/} and \code{tab/} to save graphs and tables, respectively. (We'll see why when we use Overleaf.)

\subsection{Requirements}

\begin{itemize}
  \item The R script should run without errors from the \texttt{assignment10} folder (using relative paths only)
  \item The \LaTeX{} document should compile without errors and include the table and figure produced by R
  \item Include a short \code{README.md} describing the project, what each file does, and how to reproduce the output
\end{itemize}

\newpage

% ==========================================================================
\section{Extensions}
% ==========================================================================

\subsection{Extension A: Overleaf workflow}

If you prefer to use Overleaf instead of compiling \LaTeX{} locally, adapt the project to work with Overleaf via Git:

\begin{enumerate}[label=\arabic*.]
  \item Create an Overleaf project with your \code{document.tex} file
  \item Clone the Overleaf project to your computer using Git (Overleaf provides a Git URL for each project under Menu $\rightarrow$ Git)
  \item Run your R script locally to produce the \code{.tex} table and \code{.pdf} figure
  \item Copy the output files into the cloned Overleaf folder. You can do this manually, or write a small \code{copy.sh} shell script that automates the copying. For example, this is what I use (from the Overleaf folder\footnote{Because my project folder is always in the same place (\code{Documents/Projects/}) but the Overleaf one can change. Sometimes I just clone it locally to update empirical files but work on the website if it's only about text.}):
  \begin{lstlisting}
mkdir -p tab
mkdir -p img
cp -v ~/Documents/Projects/mycurrentproject/tab/*.tex tab/
cp -v ~/Documents/Projects/mycurrentproject/img/*.pdf img/
  \end{lstlisting}
  \item Push to Overleaf via Git (\code{git add}, \code{git commit}, \code{git push}) so the document updates with the latest results
\end{enumerate}

In your \code{README.md}, describe the steps you followed and how the Overleaf integration works.

\subsection{Extension B: Makefile}

Create a \code{Makefile} in the \texttt{assignment10} folder that automates the full pipeline. Running \code{make} from the terminal should:

\begin{enumerate}[label=\arabic*.]
  \item \textbf{At least:} Run the R script to produce the table and figure
  \item \textbf{Potentially:} Compile the \LaTeX{} document to produce the final PDF
\end{enumerate}

Remember that each rule in a Makefile specifies a \textbf{target} (the output file), its \textbf{dependencies} (what it needs), and a \textbf{recipe} (the command to run). Makefile recipes must be indented with \textbf{tabs}, not spaces. \textit{(Use AI to go through this.)}

\vspace{15pt}

% Commit your work to your GitHub repository before the deadline. Everything should be inside the \texttt{assignment10} folder. Make sure your repository is public so I can access it.

% Your \texttt{assignment10} folder should contain at minimum:
% \begin{itemize}
%   \item \code{README.md} --- project description and instructions to reproduce
%   \item An R script --- analysis that produces a table and a figure
%   \item A \LaTeX{} document --- loads the R output and compiles to PDF
%   \item The dataset (\code{corruption.dta}) or instructions to obtain it
%   \item Plus any extension files (Makefile, \code{copy.sh}, etc.)
% \end{itemize}

% All R scripts should run without errors from the \texttt{assignment10} directory.

\end{document}
